\chapter{Introducción}
\label{ch.ej5:introduccion}

Este informe documenta la resolución del Ejercicio Práctico N° 5 correspondiente al Trabajo Práctico N° 2,
centrado en la generación de una señal de Modulación por Ancho de Pulso (\textit{Pulse Width Modulation} - PWM)
mediante temporizadores en el microcontrolador STM32F103C8T6 de la placa BluePill.

La técnica de PWM permite regular la cantidad de potencia entregada a una carga eléctrica (como la intensidad
lumínica de un LED o la velocidad de un motor) variando la proporción de tiempo en que la salida digital permanece
en nivel alto ($\qty{3.3}{\V}$) respecto del período total de la señal.

El objetivo de la experiencia consiste en configurar el temporizador de hardware \texttt{TIM1} para generar una
señal PWM continua e independiente sobre el canal de salida física \texttt{TIM1\_CH1} (pin \texttt{PA8}), liberando
al núcleo ARM Cortex-M3 de realizar la conmutación por software. Asimismo, mediante la variación periódica del
registro de comparación \texttt{CCR1}, se implementa una secuencia de variación gradual de brillo (\textit{fading})
en un LED exterior, manteniendo en paralelo el parpadeo determinista del LED testigo integrado (\texttt{PC13})
mediante la base de tiempo proporcionada por el temporizador del sistema (\textit{SysTick}).
